\section{The arithmetic foundation}
\label{sec:foundation}
\subsection{The theory and its finite data}
We use first-order intuitionistic logic with equality, the connectives
$\bot,\land,\lor,\to$, and the quantifiers $\forall,\exists$ over $\N$.
Negation abbreviates implication to $\bot$. The logical rules are the
ordinary introduction and elimination rules, equality reflexivity and
substitution, and the usual eigenvariable restrictions for quantifiers.
There is no rule of unrestricted excluded middle or double-negation
elimination. The nonlogical language has $0,S$, projections, and a symbol
for every function obtained by composition and primitive recursion. Its
axioms are their defining equations, $Sx\ne0$, injectivity of $S$, and
\[
 \bigl(A(0,\bar a)\land\forall n(A(n,\bar a)\to A(Sn,\bar a))\bigr)
       \to\forall n\,A(n,\bar a)
\]
for quantifier-free $A$. This is the presentation of iPRA used in the
repository \cite{proveit}. Parameters $\bar a$ remain free while induction
is applied. Universal generalization is performed afterwards.

Pairs and lists are numbers: take $\operatorname{cons}(a,b)=2^a(2b+1)$,
with zero the empty list. Head, tail, length, and indexed access are
bounded searches on the code. A signed integer is a pair $(a,b)$ read
as $a-b$; its equality to $(c,d)$ means $a+d=b+c$. A rational is a signed
numerator and a positive denominator, with equality and order determined
by cross multiplication. Coefficient lists encode polynomials in ascending
degree; matrices, vectors, permutations, and tables are lists with explicit
dimensions. Malformed input is assigned a fixed default, so that every
function is total. A theorem asserts its substantive conclusion under
the corresponding decidable well-formedness hypotheses.

\begin{lemma}[Finite arithmetic and bounded invariants]\label{found:finite}
Integer and rational arithmetic, finite polynomial arithmetic, degree-fueled
Euclidean division, finite sums and products, and Boolean tests over a
specified finite list are primitive recursive. Their usual calculation
laws and the loop invariants used below are provable with quantifier-free
induction.
\end{lemma}
\begin{proof}
Addition and multiplication satisfy their defining recursions. Distributivity
and associativity follow by induction on one numerical argument, with the
remaining arguments free; the displayed assertions are equations. Signed
arithmetic is obtained by expanding the two parts of a difference, and
rational arithmetic by multiplying positive denominators. A bounded search
keeps its current candidate and a Boolean flag saying whether a successful
entry has been found. The invariant for entries before stage $k$ is a
bounded product or sum of their characteristic functions. Thus an apparent
bounded universal quantifier is a single equation about a PR function.

Polynomial multiplication is coefficient convolution. Its associativity
reindexes finite sums, proved by induction on their bounds. Long division
subtracts a multiple of the divisor to cancel the current highest term.
The invariant $f=qg+r$, together with the decreasing degree counter, is an
equality of coefficient lists and a numerical inequality. The same argument
proves the extended Euclidean algorithm, retaining the two coefficient
vectors expressing every remainder as a combination of the inputs.
Finite permutations and tuples can be enumerated with their known length
and entry bounds, regardless of efficiency. No search in this paragraph
has a termination test without a previously supplied bound.
\end{proof}

\paragraph{A basic discipline.}
That a function is PR does not, by itself, prove every proposed property
of it in iPRA. We shall specify the numerical invariant as well. Nor is
there a PR universal evaluator for all PR programs. The algorithms below
are particular finite-data algorithms; a separately specified real name is
always a fixed PR recipe, with a numerical precision argument.

\subsection{Exact algebraic scalars without a completeness axiom}
The following construction follows the useful organization of the earlier
planar packet \cite{planar}, but all estimates needed here are stated.
The related literature on primitive recursive fields and formal algebraic
numbers supplies context, not a replacement for the bounds
\cite{selivanov,cohen}.

For $p(X)=\sum p_iX^i$ and rational $R\ge1$, define
\[
 M(p,R)=\sum_i|p_i|R^i,\qquad D(p,R)=M(p'',R).
\]
A valid scalar record is a triple $(p,l,r)$ of rational data with $l<r$,
such that, for $c=(l+r)/2$, $w=(r-l)/2$, and
$R=1+\max(|l|,|r|)$,
\begin{equation}\label{eq:record}
 p(l)p(r)<0,\qquad |p'(c)|>D(p,R)w.
\end{equation}
This is a decidable test on finite rational data. A rational $a$ is
represented by $(X-a,a-1,a+1)$.

\begin{lemma}[A reader for a scalar record]\label{found:reader}
A valid record has a PR sequence of nested rational intervals of width
at most $2^{-n}$, satisfying $p(x)=0$ at every requested precision.
There is at most one such value in the record interval.
\end{lemma}
\begin{proof}
For rational $u,v$, write
\[
 p[u,v]=\sum_{i\ge1}p_i\sum_{j<i}u^{i-1-j}v^j.
\]
The identity $p(v)-p(u)=(v-u)p[u,v]$ follows by telescoping powers.
If $u,v,c$ have absolute value at most $R$ and both $u,v$ are within
$e$ of $c$, replacement of one factor at a time gives
\begin{equation}\label{eq:divided}
 |p[u,v]-p'(c)|\le D(p,R)e.
\end{equation}
Indeed, the degree-$i$ term contributes at most
$i(i-1)|p_i|R^{i-2}e$. Thus the slope in the original interval has
a fixed sign and a positive rational lower bound.

Bisect its endpoints. If the midpoint is an exact zero, freeze at that
singleton; otherwise retain the half with opposite endpoint signs.
After $n+2+\lceil r-l\rceil$ iterations the interval width is below
$2^{-n}$. The invariant comprises rational endpoint signs, nesting,
and the width bound, so Lemma~\ref{found:finite} applies.
For opposite-sign endpoints of width $e$, the divided-difference identity
bounds each endpoint's polynomial value by $M(p',R)e$.
Consequently every desired error in $p(x)$ is obtained at an explicitly
chosen stage. If two purported roots are separated by a rational gap,
the positive slope bound forces their polynomial values to be separated
too. Refine until the polynomial errors are smaller than that separation;
this is a contradiction. No completed set of real numbers was invoked.
\end{proof}

For this argument a \emph{provisional scalar} means a particular PR
recipe for nested rational intervals with width at most $2^{-n}$.
Lemmas about readers are schemata for a fixed function $R(\bar a,n)$;
only its numerical parameters $\bar a$ and precision $n$ are inputs.
No lemma takes a code for an arbitrary PR program.
Agreement means that the two intervals meet at every precision; strict
order is witnessed by two separated intervals. Agreement is transitive:
if outer intervals have a positive gap, a sufficiently short intermediate
interval cannot meet both. The statements are formulas about numerical
precision arguments, not a new sort of infinite sequences.
Addition and multiplication have interval readers: if $|x|,|y|\le R$,
input interval widths $e$ give output widths at most $2e$ and $2Re$,
respectively. For an interval separated from zero by $b>0$, reciprocal
width is at most $e/b^2$. Intersect later intervals with the separating
one before inversion. These formulas supply an explicit backward request
for the precision needed at every input of a fixed finite expression.

\begin{lemma}[An annihilator decides sign]\label{found:sign}
Suppose a provisional scalar $x$ is accompanied by a nonzero rational
polynomial $P$ with $P(x)=0$. Its sign is determined after a PR, explicitly
bounded number of interval refinements. The zero answer means agreement
with zero at every precision, and a nonzero answer supplies a rational
separation from zero.
\end{lemma}
\begin{proof}
Remove the initial zero coefficients, writing $P=X^kG$ with $g=G(0)\ne0$.
Set
\[
 a=\min\left(\frac14,\frac{|g|}{8(1+\sum_{i>0}|G_i|)}\right)>0.
\]
For $|x|\le2a$, the nonconstant terms have total absolute value at most
$2a\sum_{i>0}|G_i|\le |g|/4$, so $G(x)$ is separated from zero.
Since $x^kG(x)=0$, this forces $x^k=0$. If $k=0$ it is impossible;
if $k>0$, an interval separating $x$ from zero by $b$ would separate
$x^k$ by $b^k$, also impossible. Hence $x=0$.

Read an interval of width less than $a$. If it lies strictly to one side
of zero, its sign and gap give the result. Otherwise it contains zero
and puts $|x|\le a$, so the preceding calculation proves the zero answer.
One may request precision $2+\lceil1/a\rceil$, since
$2^{-2-\lceil1/a\rceil}<a$. This is a bounded computation even when
the correct answer is exact equality. The proof of the zero case fixes
an arbitrary proposed separation first. For each given precision, whether
its rational interval meets zero is decidable. The disjoint alternative
supplies such a gap and is contradictory; therefore that interval meets
zero. Generalize over the precision afterwards. This is not
double-negation elimination for a universally quantified equality, and
it does not induct on a formula universally quantified over precisions.
\end{proof}

\begin{lemma}[Conversion to an isolating record]\label{found:convert}
Given the data in Lemma~\ref{found:sign}, a valid record for $x$ can be
constructed by primitive recursion.
\end{lemma}
\begin{proof}
Let $h$ be the monic squarefree part of $P$, computed by dividing out
$\gcd(P,P')$. Polynomial Euclid computes $U,V$ satisfying
$Uh+Vh'=1$. Also $P$ divides $h^{\deg P}$, so $h(x)=0$ by the
power argument in the preceding proof. These statements can be verified
without an unbounded factor search. For completeness, a bounded
factorization procedure evaluates the primitive integer associate of $P$
at $d+1$ nonzero values among $0,\ldots,2d$, where $d=\deg P$.
An integer polynomial divisor must take signed divisors of these values;
enumerate those finite choices, interpolate, and test exact divisibility.
Gauss's lemma follows by reducing a purported nonprimitive product modulo
a prime dividing all its coefficients and comparing highest nonzero
coefficients. Repeated splitting takes at most $d-1$ successful steps.
This proves the assertions about squarefree factors and their
multiplicities using finite polynomial calculations.

Take an integer $R>3+\max_{u\in I_0}|u|$, where $I_0$ is the coarse
interval of $x$, and put $B=\max(1,M(V,R))$ and
$L=1+D(h,R)$. The Bezout identity at $x$ gives
$|h'(x)|\ge1/B$. Choose $\delta=2^{-s}<1/(32BL)$ and read an interval
of width at most $\delta$ with midpoint $c$. The slope estimate
\eqref{eq:divided} shows that throughout $[c-2\delta,c+2\delta]$
the divided differences have the sign of $h'(x)$ and absolute value
greater than $1/(2B)$. Its two endpoints are on opposite sides of $x$
at distance at least $\delta$. Their values under $h$ therefore have
opposite signs, and the second test in \eqref{eq:record} also holds.
Its own test radius $1+\max|c\pm2\delta|$ is smaller than $R$, so
its derivative majorant is no larger than $D(h,R)$.
The triple $(h,c-2\delta,c+2\delta)$ is the required record; uniqueness
in Lemma~\ref{found:reader} identifies its reader with $x$.
All precisions are explicit rational functions of finite coefficient bounds.
\end{proof}

\begin{theorem}[A coded ordered field]\label{found:field}
Valid records, modulo agreement, carry PR addition, subtraction,
multiplication, division with $0^{-1}=0$, and decidable equality and order.
All ordered-field laws, including inverse laws under a nonzero hypothesis,
hold in the numerical interpretation in iPRA. Write $\K$ for this field
notation; internally no quotient type is added.
\end{theorem}
\begin{proof}
For monic annihilators $p_1,p_2$ of degrees $d_1,d_2$, form the finite
rational algebra
\[
 E=\Q[T_1,T_2]/(p_1(T_1),p_2(T_2)).
\]
The monomials $T_1^iT_2^j$, $i<d_1,j<d_2$, are a coefficient basis.
Reduction in the two separate monic polynomials gives an explicit
multiplication table. For $v=T_1+T_2,T_1-T_2,T_1T_2$ compute the matrix
of multiplication by $v$. Its characteristic polynomial is a nonzero
monic annihilator. To see the annihilation directly, expand
$(YI-A)\operatorname{adj}(YI-A)=\det(YI-A)I$, compare coefficients,
multiply by successive powers of $A$, and add. All intermediate matrix
terms telescope, leaving the characteristic polynomial at $A$ equal to
zero. Applying it to the vector $1$ gives the required relation for $v$.
Evaluation at the input scalars preserves the reductions even if the
two inputs satisfy additional algebraic relations.

Use interval arithmetic for the provisional result and
Lemma~\ref{found:convert} for its output record. Negation is similar.
For an input whose sign is nonzero, its rational gap licenses interval
reciprocal, with annihilator $X^{\deg p}p(1/X)$. The zero branch returns
zero. Compare two scalars by applying Lemma~\ref{found:sign} to their
difference. Different records may represent the same scalar.
This test returns zero exactly when the readers agree: if the two
intervals meet at every equal precision, their interval differences contain
zero at every refinement. Conversion preserves their provisional value,
so its sign is zero. Conversely, disjoint input intervals supply a positive
rational gap that persists under refinement; their difference cannot have
zero sign. This proves that the PR equality used in subsequent formulas
coincides with the reader agreement relation.

The field identities are rational identities evaluated at arbitrary
precision. Given a claimed positive output gap, propagate the requested
error backwards through their finite expression trees and use the rational
identity on sufficiently accurate input approximations. The error becomes
smaller than the gap, disproving the alleged separation. The inverse
identity uses the supplied input gap. Translation preserves a positive
gap after refining the common summand, and multiplication preserves
positivity because positive lower bounds $b,c$ give lower bound $bc$.
Trichotomy is exactly the three finite branches of the sign test.
Thus both arithmetic and its order laws are proved, rather than assumed.
\end{proof}

\begin{theorem}[Polynomial roots, square roots, and finite root lists]
\label{found:roots}
Over $\K$, a polynomial with opposite nonzero signs at two endpoints
has a constructible-in-arithmetic zero strictly between them. Every
nonnegative scalar has a unique nonnegative square root. Every
nonzero polynomial has a PR list of all its distinct real roots,
in increasing order. These statements have iPRA proofs in the coded
interpretation.
\end{theorem}
\begin{proof}
Normalize a nonconstant polynomial
$f(Z)=Z^n+\sum_{i<n}x_iZ^i$ to be monic. Form the finite rational
algebra on one generator $T_i$ for every coefficient record, modulo its
monic annihilator, and adjoin $Z$ modulo the monic polynomial
$Z^n+\sum T_iZ^i$. Multiplication by $Z$ now has a rational matrix.
Its characteristic polynomial $H$ is monic and satisfies
\begin{equation}\label{eq:eliminant}
 f(y)=0\quad\Longrightarrow\quad H(y)=0.
\end{equation}
This follows from the same adjugate identity and evaluation argument as
before. The matrix dimension is $n\prod_i d_i$, a known natural number.

For endpoints $s<t$ with opposite signs, substitute $s+(t-s)T$ and
bisect the rational interval $[0,1]$. Exact signs are already available
from Theorem~\ref{found:field}; they do not call the root finder being
defined. After $k$ steps the width is $2^{-k}$. A rational upper bound
on the derivative's coefficients bounds the polynomial residual by
the interval width times that constant. Thus the provisional value $\tau$
is a zero of the substituted polynomial. Put $y=s+(t-s)\tau$ using its
fixed interval readers. Then $f(y)=0$, so
relation~\eqref{eq:eliminant} for $f$ and the conversion lemma return
a finite algebraic record for $y$. Nonzero endpoint signs exclude either endpoint.
For $b>0$, apply this to $T^2-b$ between $0$ and $b+1$; for $b=0$
return zero. If $r,s\ge0$ and $r^2=s^2$, then
$(r-s)(r+s)=0$ and either both vanish or $r+s>0$, proving uniqueness.

For exhaustive root lists take the squarefree part $h$ of $H$ and
an integer $R>2+\sum_{i<\deg h}|h_i|$. Leading-term dominance puts
every root of $h$ strictly inside $[-R,R]$. Compute $U,V$ with
$Uh+Vh'=1$, and $B,L$ as above at radius $R+1$. Partition $[-R,R]$
into rational dyadic cells of length $\delta<1/(32BL)$. Record every
grid zero and every cell whose endpoint signs differ, obtaining its
zero by bisection and conversion. Any root $z$ lies either on the grid
or in a cell. The Bezout derivative bound at $z$ makes the slope
strictly signed throughout that cell, so an interior root forces an
endpoint sign change and is unique there. A grid root similarly excludes
an extra root in an adjacent cell. Filter the finite result by the exact
test $f(z)=0$. Relation~\eqref{eq:eliminant} proves that every desired
root survives. Constants and the zero polynomial are handled separately,
the latter by a flag rather than an impossible finite list of all scalars.

The grid length, degrees, coefficient lists, and bisection precision are
computed before their loops. The grid invariant is the bounded Boolean
fold of processed sign changes and the stored root records. Completeness
is proved for an arbitrary candidate root as a free parameter. It never
uses induction on an assertion quantified over all roots or all precisions.
\end{proof}

Externally, the records denote exactly the real algebraic numbers.
Every record denotes its real root by the nested-interval construction,
and the exhaustive procedure finds every real root of a rational
polynomial. This external identification explains the geometric meaning.
The arithmetic proofs themselves use only the finite definitions above.

\subsection{A precision-level version that is visibly first order}
For nonnegative rational $r$ and $\epsilon=2^{-n}$, let $m$ be the largest
integer in the bounded range $0\le m\le 2^n(1+\lceil r\rceil)$ with
$m^2\le 2^{2n}r$. Cross multiplication removes the rational denominator.
Then
\begin{equation}\label{eq:meshroot}
 m\epsilon\le\sqrt r<(m+1)\epsilon,
 \qquad 0\le r-(m\epsilon)^2<(2m+1)\epsilon^2.
\end{equation}
The upper bound follows because $m+1$ fails the inequality; existence of
the bounded maximum follows by its Boolean scan. This is an explicit
quantifier-free arithmetic proof of a root approximation, not an appeal
to a limit inside the object theory.

For later reference the elementary inequalities
\begin{equation}\label{eq:sqrtholder}
 |\sqrt r-\sqrt s|^2\le |r-s|\quad(r,s\ge0),\qquad
 |r^{-1}-s^{-1}|\le\frac{|r-s|}{b^2}\quad(r,s\ge b>0)
\end{equation}
follow by factoring the difference of squares and the difference of
reciprocals. They give a concrete path from every fixed finite construction
to PR error bounds for its rational observations.
